\paragraph{Workload definition.}
We consider a complex information need $q$ over one or more evidence sources $\mathcal{D}$ under a finite global retrieval-call budget $B$. Solving $q$ requires acquiring a set of mutually dependent evidence items before producing an answer (or, for verification tasks, a verdict). Two properties make this problem non-trivial. First, the evidence items are \emph{dependent}: binding one item (e.g., an entity) is a precondition for acquiring another (e.g., a relation that connects it to a second entity). Second, the system operates under \emph{partial information}: whether a bridge binding exists, how many candidates a retrieval step will return, and whether the answer model can use the materialized evidence are all uncertain before execution. A plan must therefore commit resources before the outcomes that determine its success are observed.

\paragraph{Typed evidence plan.}
A typed evidence plan is a tuple $P = (S, J, O)$ where $S$ is a set of \emph{evidence slots} (typed requirements over what evidence must be acquired), $J$ is a set of \emph{join edges} (specifying which slots must share bound variables), and $O$ is a set of \emph{relational operators} (compare, extremum, polar) over slot outputs. Each slot $s \in S$ carries a type (entity, passage, relation), a set of variables to bind, and a status (unresolved, partial, satisfied). Join edges encode structural coupling: slot $s_j$ depends on bindings produced by slot $s_i$ iff $(s_i, s_j) \in J$.

\paragraph{Structural evidence graph.}
The join edges $J$ induce a directed acyclic graph $G_P = (S, J)$ over the slots. The structural depth of $P$, denoted $\texttt{structural\_hops}(P)$, is the length of the longest path in $G_P$ measured in join and operator edges. This is a deterministic property of the plan's topology, computed without any LLM call. Plans with $\texttt{structural\_hops} \geq 2$ involve at least one intermediate binding through which evidence must flow; plans with $\texttt{structural\_hops} < 2$ have only shallow or independent slot structure.

\paragraph{Physical allocation.}
A physical allocation $A(P) = \{a_s \mid s \in S\}$ assigns a non-negative integer number of retrieval calls $a_s$ to each slot. Under the \emph{static} allocation strategy, each slot $s$ receives a fixed local budget: $a_s = \lfloor B / |S| \rfloor + \mathbb{1}_{s \leq B \bmod |S|}$, distributing calls as evenly as possible. Under the \emph{budget-aware flat} strategy, allocations are globally optimized to maximize evidence acquisition under the constraint $\sum_s a_s \leq B$, with all slots treated as equally important.

\paragraph{Allocation feasibility.}
Define the allocation feasibility predicate: $F(P, A, B) = \mathbf{1}\!\left[\sum_{s \in S} a_s \leq B\right]$. When $F = 0$, the allocation exceeds the global budget: the static executor must interrupt materialization before completing all slots, resulting in a \emph{budget-exhaustion event} (BE). When $F = 1$, the allocation fits within the budget; whether the plan succeeds depends on retrieval quality and answer generation, not on budget exhaustion.

\paragraph{Budget-exhaustion as a necessary condition.}
Under the static allocation implementation, every observed budget-exhaustion event satisfies $F = 0$: the sum of static allocations exceeded the global budget. Formally, $\Pr[F = 0 \mid \text{BE observed}] \geq \Pr[F = 0 \mid \text{BE observed}]$. The contrapositive is $F = 1 \Rightarrow \text{no BE under static execution}$. This is a necessary condition, not a sufficient one: a plan with $F = 0$ may still terminate without budget exhaustion if the executor's other stopping criteria (step limit, LLM-call limit, early convergence) trigger first.

\paragraph{Objective.}
We formulate physical planning as constrained optimization: maximize expected evidence acquisition and answer quality subject to the global retrieval-call budget $B$. The key insight is that the logical plan $P$ is fixed (determined at compile time by the slot compiler), while the physical allocation $A$ is the optimization variable. The budget-feasibility diagnosis operates on $P$ at compile time; the dispatch policy selects the allocation strategy at execution time based on structural depth.
